\bookchapter{第四十七章：速查手冊同資源}

\chapterintro{本章係一份可以獨立使用嘅速查手冊。當你記唔清某個原則或者需要快速決策嗰陣，直接揭到嚟呢度。建議你將佢當係「口袋卡」：五大原則速查、工具決策樹、常見問題、學習資源，一頁都唔漏。}

\booksection{一、思維模型速查卡}

\begin{center}
\begin{tabularx}{\linewidth}{|X|X|X|}
\hline
\textbf{原則} & \textbf{一句講曬} & \textbf{操作要點} \tabularnewline
\hline
Think Before Coding & 先諗清楚先開口 & 答三條問題：做乜嘢？結果係點嘅？有乜嘢約束？ \tabularnewline
Simplicity First & 用最少代碼解決問題 & 每次只做一個功能，斬走唔必要嘅複雜度 \tabularnewline
Surgical Changes & 精準修改唔越界 & 明確指定修改範圍，拒絕AI嘅「順手改嚇」 \tabularnewline
Goal-Driven Execution & 講「做乜嘢」唔好講「點樣做」 & 畀目標唔畀步驟，畀AI自主揀路徑 \tabularnewline
Iterate \& Verify & 產生-驗證-修正循環 & 三層驗證：功能、邊界、架構 \tabularnewline
\hline
\end{tabularx}
\end{center}

\begin{center}\small\color{gray}五大思維原則速查卡\end{center}

\booksection{二、工具選擇決策樹}

當你需要揀 AI 編程工具嗰陣，跟以下決策路徑嚟揀。留意：呢個唔係「邊個最好」嘅排名，而係「邊個最啱你」嘅判斷框架。

你想做乜嘢？

零基礎，想快速體驗 揀雲端 Builder (v0/Bolt/Replit)

日常開發，需要代碼補全 揀 IDE 內置 AI (Cursor/Trae)

大型項目，需要批量操作 揀終端 Agent (Claude Code)

快速原型驗證想法 揀雲端 Builder 或 Trae SOLO

\begin{center}\small\color{gray}圖 6：AI 編程工具選擇決策樹\end{center}

\booksection{三、常見問題 FAQ}

Q1：我完全唔識編程，可以用 AI 編程嗎？

得。2026 年嘅 AI 編程工具已經夠曬強大，你完全可以用自然語言描述需求，叫 AI 產生代碼。但有兩點建議：第一，由雲端 Builder 工具（例如 v0、Bolt）開始，佢哋唔使配置開發環境；第二，邊用邊學基礎概念——唔使精通語法，但你需要可以「睇得明」AI 寫嘅代碼嘅大意，咁樣先可以有效驗收。

Q2：我應該學邊個編程語言？

睇你想做乜嘢。做網站前端——JavaScript/TypeScript；做數據處理或者 AI——Python；做企業級後端——Java 或者 Go。但更重要嘅係：喺 AI 編程時代，「語言」嘅重要性喺度下降。AI 可以幫你處理語法細節，你需要明白嘅係「程序嘅結構」——乜嘢係函數、乜嘢係循環、乜嘢係數據結構。呢啲概念係跨語言嘅，一旦明白咗，轉乜嘢語言都快。

Q3：AI 寫嘅代碼有冇版權問題？

截至 2026 年 8 月，呢個仍然係一個有爭議嘅法律問題。一般嚟講，AI 產生嘅代碼喺大部分司法管轄區唔會被視為有明確嘅版權擁有者。但如果你用喺商業項目，建議：(1) 瞭解你用緊嘅 AI 工具嘅服務條款入面關於知識產權嘅條款；(2) 對關鍵代碼做人工審查同修改，令佢有你嘅創造性投入；(3) 避免叫 AI 直接抄可能有版權嘅開源代碼。

Q4：AI 會唔會取代程序員？

短期內唔會。AI 取代嘅唔係「程序員」呢個職業，而係「逐行編碼」呢樣工作。就好似計數機冇取代數學家，而係令數學家由繁瑣嘅手算入面解放出嚟咁。AI 編程工具令程序員由「寫代碼」轉向「諗問題」——呢樣嘢實際上提高咗程序員嘅價值上限。將來嘅競爭力唔喺於「可以寫幾多行代碼」，而喺於「可以解決幾複雜嘅問題」。

Q5：點樣判斷 AI 嘅產出啱唔啱？

三層驗證法。第一層：功能驗證——佢行唔行到？輸入輸出啱唔啱？呢個係最基本嘅。第二層：邊界驗證——空輸入、超長輸入、異常輸入會點？呢層可以發現隱藏嘅問題。第三層：架構驗證——代碼結構合唔合理？有冇寫死嘅嘢？方唔方便擴展？呢層需要你有一定嘅技術判斷力。如果你判斷唔到第三層，至少做到頭兩層——行到、邊界正常，就已經比大部分初學者強。

\booksection{四、推薦學習資源}

以下係截至 2026 年 8 月嘅推薦學習資源。留意，互聯網資源連結可能會變，如果某個連結失效，搜尋佢個標題通常可以搵到最新版本。

\begin{center}
\begin{tabularx}{\linewidth}{|X|X|X|}
\hline
\textbf{類型} & \textbf{資源} & \textbf{推薦理由} \tabularnewline
\hline
原始文獻 & Andrej Karpathy 嘅 CLAUDE.md & 五大原則嘅源頭，70行但字字珠璣 \tabularnewline
原始文獻 & Anthropic 官方文檔 (Claude Code) & 終端Agent嘅最佳實踐指南 \tabularnewline
教程 & Vibe Coding 社羣整理嘅 Skills 合集 & 社羣實踐者總結嘅實戰技巧集合 \tabularnewline
教程 & Cursor 官方文檔同 YouTube 頻道 & IDE內置AI嘅權威使用指南 \tabularnewline
社羣 & GitHub 上面嘅 awesome-claude-code & Claude Code 嘅最佳實踐同插件集 \tabularnewline
社羣 & Reddit 嘅 r/cursor 社羣 & 全球Cursor用戶嘅經驗分享同討論 \tabularnewline
基礎 & MDN Web Docs (developer.mozilla.org) & 前端開發最權威嘅參考文檔 \tabularnewline
基礎 & Python 官方教程 (docs.python.org) & Python學習嘅官方起點 \tabularnewline
\hline
\end{tabularx}
\end{center}

\begin{center}\small\color{gray}推薦學習資源清單\end{center}

\booksection{五、結語：工具喺手，方法論喺心}

返返去本書嘅核心信念：工具每三個月迭代一次，方法論每三年先進化一次。喺 2026 年 8 月嘅今日，AI 編程工具嘅能力已經強大到令人驚嘆——你可以用一句話整一個網站，用一段描述產生一個全端應用，用一條指令完成一個複雜重構。但工具越強大，方法論嘅重要性反而越突出——因為工具嘅強大會放大你嘅每一個決策。

如果你嘅需求清晰，AI 可以幫你高效實現；如果你嘅需求含糊，AI 都會高效咁做出錯嘅嘢。如果你跟簡單至上，AI 可以幫你快速驗證諗法；如果你貪多求全，AI 都會幫你快速堆積技術債。如果你精準修改，AI 可以幫你好似外科手術咁優化代碼；如果你指令含糊，AI 都會幫你大刀闊斧咁引入 bug。

呢個就係點解本書花咗大部分篇幅講方法論而唔係工具操作。唔係因為工具唔重要，而係因為方法論係工具嘅「乘數」——方法論越好，工具嘅價值越大；方法論係零，工具幾強都只係零乘任何數。

你唔再係一個人喺度寫程式。你係一個指揮官，AI 係你嘅軍團。你思考嘅質素決定咗軍團嘅戰鬥力。由今日開始，用方法論武裝自己嘅思維，用工具放大自己嘅能力——呢個就係 AI 編程嘅全部祕密。

\booksection{參考文獻}

\begin{enumerate}
  \item Andrej Karpathy. (2025). Vibe Coding. Social media post introducing the concept of vibe-driven programming with AI.
  \item Andrej Karpathy. (2025). CLAUDE.md - Project guidelines for AI-assisted development. GitHub repository, 70+K stars.
  \item Anthropic. (2026). Claude Code: Terminal-native AI programming. Official documentation and best practices.
  \item Cursor Team. (2026). Cursor 3: The next generation of AI-first IDE. Release notes and documentation.
  \item OpenAI. (2026). Codex CLI: Bringing AI to the terminal. Technical documentation.
  \item Model Context Protocol (MCP) Specification. (2025-2026). Open protocol for AI-tool integration.
  \item Vibe Coding Community. (2026). Community-curated skills and best practices for AI-assisted programming.
  \item Spec-Driven Development (SDD) methodology. (2026). Emerging practices in specification-first AI programming.
\end{enumerate}
